% Phương pháp luận (Chương 2) và Thực nghiệm (Chương 3) của DyTimeBase, viết theo form Đồ án I.
% Biên dịch thử từ thư mục docs/: pdflatex dytimebase_phuong_phap_luan.tex (chạy 2 lần).
% Các bảng ở docs/tables/ sinh tự động từ results/*.json bằng: python scripts/make_report_tables.py
\documentclass[12pt,a4paper]{report}
\usepackage[utf8]{inputenc}
\usepackage[T5]{fontenc}
\usepackage[a4paper,margin=2.5cm]{geometry}
\usepackage{amsmath,amssymb,bm}
\usepackage{booktabs,array,graphicx,float,multirow,rotating}
\renewcommand{\chaptername}{Chương}
\renewcommand{\tablename}{Bảng}
\renewcommand{\figurename}{Hình}
\renewcommand{\bibname}{Tài liệu tham khảo}
\setlength{\emergencystretch}{2em}
\numberwithin{equation}{chapter}

\newcommand{\R}{\mathbb{R}}
\newcommand{\one}{\mathbf{1}}
\newcommand{\LN}{\operatorname{LayerNorm}}
\newcommand{\GELU}{\operatorname{GELU}}
\newcommand{\DF}{\operatorname{DF}}
\newcommand{\FFN}{\operatorname{FFN}}
\newcommand{\vecr}{\operatorname{vec}}

\begin{document}
\setcounter{chapter}{1}
\chapter{Phương pháp luận}

Chương này trình bày nền tảng toán học của bài toán dự báo chuỗi thời gian đa biến và cấu trúc của mô hình
\textbf{DyTimeBase}. DyTimeBase giữ nguyên xương sống tuyến tính, độc lập kênh của TimeBase (rút một số ít
\emph{chuỗi cơ sở} từ các đoạn theo chu kỳ, rồi dự báo theo đoạn), và bổ sung một nhánh \emph{trộn kênh} lấy từ
DyTimeNet (khối CrossVBlock với chiến lược DropFusion và khối IntraVBlock). Điểm khác biệt cốt lõi là nhánh trộn kênh
hoạt động trên \emph{không gian cơ sở} gồm $P\cdot R$ số mỗi kênh, thay vì trên toàn bộ $L$ điểm của chuỗi gốc.

%=====================================================================
\section{Mô hình hóa toán học}

\subsection{Dữ liệu đầu vào}

Xét một hệ thống quan sát gồm $C$ biến số thay đổi theo thời gian. Một mẫu đầu vào là cửa sổ lịch sử độ dài $L$:
\begin{equation}
\mathcal{X}_{train} =
\begin{bmatrix}
x^{(1)}_1 & x^{(2)}_1 & \dots & x^{(C)}_1\\
\vdots & \vdots & \ddots & \vdots\\
x^{(1)}_L & x^{(2)}_L & \dots & x^{(C)}_L
\end{bmatrix} \in \R^{L\times C},
\qquad \mathbf{x}^{(i)} = \big(x^{(i)}_1,\dots,x^{(i)}_L\big)^{\!\top}\in\R^{L},
\label{eq:input}
\end{equation}
trong đó $x^{(i)}_t$ là giá trị của biến $X^{(i)}$ tại thời điểm $t$, đã được chuẩn hóa Z-score theo từng biến với
trung bình và độ lệch chuẩn tính trên tập huấn luyện. Nhãn tương ứng là $S$ bước tương lai:
\begin{equation}
\mathcal{X}_{test} =
\begin{bmatrix}
x^{(1)}_{L+1} & \dots & x^{(C)}_{L+1}\\
\vdots & \ddots & \vdots\\
x^{(1)}_{L+S} & \dots & x^{(C)}_{L+S}
\end{bmatrix} \in \R^{S\times C}.
\label{eq:label}
\end{equation}

Ngoài $L$, $S$, $C$, mô hình dùng thêm các kích thước sau:
\begin{itemize}
  \item $P$: độ dài đoạn, bằng chu kỳ tự nhiên của dữ liệu (ví dụ $P=24$ với dữ liệu theo giờ có chu kỳ ngày).
  \item $N=\lceil L/P\rceil$ và $N'=\lceil S/P\rceil$: số đoạn của quá khứ và của tương lai.
  \item $R\ll N$: số chuỗi cơ sở (basis).
  \item $D$: số chiều ẩn của nhánh trộn kênh; $C'=\lfloor 0.75\,C\rfloor$: số chiều nút thắt của CrossVBlock;
        $D_{ff}=2D$: số chiều ẩn của FFN trong IntraVBlock.
\end{itemize}

\subsection{Đầu ra và Hàm mục tiêu}

Mô hình $\mathcal{F}$ (với tập tham số $\Theta$) nhận $\mathcal{X}_{train}$ và ước lượng đồng thời
$S$ bước tương lai của cả $C$ biến:
\begin{equation}
\hat{\mathcal{X}} = \mathcal{F}(\mathcal{X}_{train}) \in \R^{S\times C}.
\label{eq:model}
\end{equation}

Sai số dự báo được đo bằng trung bình bình phương sai số:
\begin{equation}
\mathcal{L}_{MSE} = \frac{1}{S\times C}\sum_{t=1}^{S}\sum_{i=1}^{C}\Big(x^{(i)}_{L+t}-\hat{x}^{(i)}_{L+t}\Big)^2 .
\label{eq:mse}
\end{equation}

Khác với DyTimeNet, hàm mục tiêu khi huấn luyện DyTimeBase có thêm một số hạng điều chuẩn trực giao
$\mathcal{L}_{orth}$ trên các chuỗi cơ sở. Số hạng này được định nghĩa ở mục~\ref{sec:basis}, công
thức~\eqref{eq:orth}:
\begin{equation}
\mathcal{L} = \mathcal{L}_{MSE} + \lambda_{orth}\,\mathcal{L}_{orth},
\label{eq:objective}
\end{equation}
trong đó $\lambda_{orth}\ge 0$ là siêu tham số. $\mathcal{L}_{orth}$ chỉ dùng khi huấn luyện; khi đánh giá, mô hình
chỉ được đo bằng MSE và MAE.

%=====================================================================
\section{Tổng quan kiến trúc của mô hình}

Hình~\ref{fig:dytimebase} mô tả luồng dữ liệu của DyTimeBase. Mỗi mẫu đi qua các khối sau:
\begin{enumerate}
  \item \textbf{Phân đoạn}: chuỗi độ dài $L$ của mỗi kênh được cắt thành $N$ đoạn dài $P$, xếp thành ma trận
        $P\times N$ (hàng là pha trong chu kỳ, cột là đoạn).
  \item \textbf{Chuẩn hóa theo pha}: trừ trung bình của từng pha, tính trên $N$ đoạn.
  \item \textbf{Trích xuất cơ sở}: một ánh xạ tuyến tính $N\to R$ nén $N$ đoạn thành $R$ chuỗi cơ sở,
        cho ma trận $\mathbf{E}^{(i)}\in\R^{P\times R}$.
  \item \textbf{Trộn kênh trong không gian cơ sở}, gồm bốn bước:
        mã hóa $P\cdot R\to D$, CrossVBlock (trộn $C$ kênh với DropFusion), IntraVBlock (FFN theo từng kênh), và chiếu
        ngược $D\to P\cdot R$ (khởi tạo bằng 0) để ra phần hiệu chỉnh $\Delta\mathbf{E}^{(i)}$.
  \item \textbf{Dự báo theo đoạn}: một ánh xạ tuyến tính $R\to N'$ trên
        $\mathbf{E}'^{(i)}=\mathbf{E}^{(i)}+\Delta\mathbf{E}^{(i)}$.
  \item \textbf{Khôi phục và ghép}: cộng lại trung bình theo pha, trải $N'$ đoạn thành chuỗi và lấy $S$ bước đầu.
\end{enumerate}

\begin{figure}[ht]
  \centering
  % \includegraphics[width=\textwidth]{dytimebase.png}
  \fbox{\parbox{0.9\textwidth}{\centering\vspace{2cm}(Sơ đồ kiến trúc DyTimeBase)\vspace{2cm}}}
  \caption{Overview Pipeline của mô hình DyTimeBase}
  \label{fig:dytimebase}
\end{figure}

Mọi phép biến đổi theo trục thời gian đều chạy trên các kích thước $N$, $R$, $P$ (cỡ vài chục) thay vì $L$, $S$
(cỡ hàng trăm). Phần trộn kênh chạy trên vector $D$ chiều của mỗi kênh. Vì vậy chi phí tính toán thấp hơn nhiều
so với DyTimeNet (mục~\ref{sec:complexity}).

%=====================================================================
\section{Xây dựng mô hình DyTimeBase}

Các khối từ~\ref{sec:segment} đến~\ref{sec:basis} và khối~\ref{sec:forecast} áp dụng \emph{độc lập cho từng kênh}
$i=1,\dots,C$. Chỉ khối~\ref{sec:mixing} trao đổi thông tin giữa các kênh.

\subsection{Phân đoạn theo chu kỳ và Chuẩn hóa theo pha}
\label{sec:segment}

\textbf{a) Phân đoạn.} Khi $L$ không chia hết cho $P$, chuỗi được đệm thêm $\delta=NP-L$ giá trị ở đầu, lấy từ đoạn
đầu tiên:
\begin{equation}
\bar{\mathbf{x}}^{(i)} = \big(x^{(i)}_{P-\delta+1},\dots,x^{(i)}_{P},\;x^{(i)}_1,\dots,x^{(i)}_L\big)^{\!\top}\in\R^{NP}
\label{eq:pad}
\end{equation}
($\delta=0$ thì $\bar{\mathbf{x}}^{(i)}=\mathbf{x}^{(i)}$). Chuỗi sau đó được xếp thành ma trận đoạn
$\mathbf{X}^{(i)}\in\R^{P\times N}$, mỗi cột là một đoạn dài $P$:
\begin{equation}
\big[\mathbf{X}^{(i)}\big]_{p,n} = \bar{x}^{(i)}_{(n-1)P+p},\qquad p=1,\dots,P,\; n=1,\dots,N.
\label{eq:segment}
\end{equation}
Hàng $p$ của $\mathbf{X}^{(i)}$ gom các giá trị cùng pha $p$ trong chu kỳ (ví dụ cùng giờ trong ngày) của $N$ chu kỳ
liên tiếp.

\textbf{b) Chuẩn hóa theo pha (Period Normalization).} Mỗi pha có mức trung bình riêng (ví dụ ban ngày phụ tải
cao hơn ban đêm), nên mô hình trừ trung bình của từng pha:
\begin{equation}
\bm{\mu}^{(i)} = \frac{1}{N}\,\mathbf{X}^{(i)}\one_N \in\R^{P},\qquad
\tilde{\mathbf{X}}^{(i)} = \mathbf{X}^{(i)} - \bm{\mu}^{(i)}\one_N^{\top}\in\R^{P\times N}.
\label{eq:period-norm}
\end{equation}
\textbf{Tác dụng và ý nghĩa:} Phép trừ này tách mẫu hình lặp lại trong một chu kỳ ra khỏi chuỗi, để các khối phía sau
chỉ phải học phần dao động quanh mẫu hình đó. Bước này đóng vai trò tương tự RevIN trong DyTimeNet nhưng không chia
cho độ lệch chuẩn. $\bm{\mu}^{(i)}$ được giữ lại để cộng trả ở bước cuối~\eqref{eq:denorm}.

\subsection{Khối Trích xuất cơ sở (Basis Extraction)}
\label{sec:basis}

Một ánh xạ tuyến tính dùng chung cho mọi kênh nén $N$ đoạn thành $R$ chuỗi cơ sở:
\begin{equation}
\mathbf{E}^{(i)} = \tilde{\mathbf{X}}^{(i)}\mathbf{W}_B + \one_P\mathbf{b}_B^{\top}\in\R^{P\times R},
\qquad \mathbf{W}_B\in\R^{N\times R},\;\mathbf{b}_B\in\R^{R}.
\label{eq:basis}
\end{equation}
Cột thứ $r$ của $\mathbf{E}^{(i)}$ là $\sum_{n=1}^{N}[\mathbf{W}_B]_{n,r}\,\tilde{\mathbf{X}}^{(i)}_{:,n}+[\mathbf{b}_B]_r$,
tức một tổ hợp tuyến tính của $N$ đoạn quá khứ. Mỗi cột vì thế là một ``đoạn đại diện'' độ dài $P$, và cả cửa sổ
$L$ điểm được tóm tắt bằng $R$ đoạn đại diện như vậy.

\textbf{Ràng buộc trực giao.} Để $R$ chuỗi cơ sở mang thông tin khác nhau và không trùng lặp, mô hình phạt các tích
vô hướng giữa hai chuỗi cơ sở khác nhau. Gọi $\mathbf{G}^{(i)}=\mathbf{E}^{(i)\top}\mathbf{E}^{(i)}\in\R^{R\times R}$
là ma trận Gram. Với một lô gồm $B$ mẫu:
\begin{equation}
\mathcal{L}_{orth} = \frac{1}{B\,C}\sum_{b=1}^{B}\sum_{i=1}^{C}
\Big\|\mathbf{G}^{(i)}_b - \operatorname{diag}\big(\mathbf{G}^{(i)}_b\big)\Big\|_F ,
\label{eq:orth}
\end{equation}
trong đó $\operatorname{diag}(\cdot)$ giữ lại đường chéo và $\|\cdot\|_F$ là chuẩn Frobenius. Số hạng này đạt giá
trị nhỏ nhất bằng 0 khi các cột của $\mathbf{E}^{(i)}$ đôi một trực giao.

\subsection{Khối Trộn kênh trong không gian cơ sở}
\label{sec:mixing}

Các khối trước xử lý từng kênh riêng rẽ. Khối này cho các kênh trao đổi thông tin. Chỉ hai khối CrossVBlock (cùng
DropFusion) và IntraVBlock được giữ nguyên từ DyTimeNet. Hai lớp bao quanh chúng là thiết kế mới của DyTimeBase:
\begin{itemize}
  \item Lớp mã hóa $PR\to D$ thay cho phép chiếu thời gian $L\to D$ của DyTimeNet. Nó nhận các chuỗi cơ sở
        $\mathbf{E}^{(i)}$ ($P\cdot R$ số mỗi kênh) thay vì chuỗi gốc ($L$ số mỗi kênh).
  \item Lớp chiếu ngược $D\to PR$, khởi tạo bằng 0, thay cho lớp chiếu ra $D\to S$ của DyTimeNet.
\end{itemize}

\textbf{a) Làm phẳng và Mã hóa.} Mỗi ma trận cơ sở được làm phẳng theo hàng,
$\mathbf{e}^{(i)}=\vecr(\mathbf{E}^{(i)})\in\R^{PR}$ với $[\mathbf{e}^{(i)}]_{(p-1)R+r}=[\mathbf{E}^{(i)}]_{p,r}$.
Các vector này được ghép thành $\bar{\mathbf{E}}=[\mathbf{e}^{(1)},\dots,\mathbf{e}^{(C)}]\in\R^{PR\times C}$ rồi chiếu
xuống không gian ẩn $D$ chiều. Phép biến đổi này cùng dạng với phép mã hóa thời gian
$\mathcal{X}_{enc}=\mathbf{W}_{enc}^{\top}\mathcal{X}'+\mathbf{b}_{enc}$ của DyTimeNet (Đồ án I, công thức (2.6)), nhưng
đầu vào là $PR$ hệ số cơ sở thay cho $L$ điểm:
\begin{equation}
\mathcal{X}_{enc} = \mathbf{W}_{enc}^{\top}\bar{\mathbf{E}} + \mathbf{b}_{enc}\one_C^{\top}\in\R^{D\times C},
\qquad \mathbf{W}_{enc}\in\R^{PR\times D},\;\mathbf{b}_{enc}\in\R^{D}.
\label{eq:embed}
\end{equation}

\textbf{b) Khối Liên biến CrossVBlock với DropFusion.} Đầu tiên, $\mathcal{X}_{enc}$ được chuẩn hóa LayerNorm theo
từng cột (từng kênh, trên $D$ đặc trưng):
\begin{equation}
\mathcal{X}'_{enc} = \LN(\mathcal{X}_{enc})\in\R^{D\times C}.
\label{eq:ln-cross}
\end{equation}
\textbf{Tác dụng và ý nghĩa:} LayerNorm đưa biểu diễn của mọi kênh về cùng thang đo, để phép trộn kênh phía sau không
bị chi phối bởi các kênh có biên độ lớn.

Phép trộn kênh dùng hàm kích hoạt GELU:
\begin{equation}
\GELU(x) = x\cdot\Phi(x) = \frac{x}{2}\left(1+\operatorname{erf}\left(\frac{x}{\sqrt{2}}\right)\right),
\label{eq:gelu}
\end{equation}
với $\Phi(x)$ là hàm phân phối tích lũy của phân phối chuẩn $\mathcal{N}(0,1)$.

Mỗi lớp tuyến tính trên trục kênh được thay bằng lớp DropFusion. Với đầu vào $\mathbf{Z}\in\R^{D\times C_{in}}$ và tham
số $\mathbf{W}_{df}\in\R^{C_{in}\times C_{out}}$, $\mathbf{b}_{df}\in\R^{C_{out}}$, đầu tiên mặt nạ cạnh được áp lên
ma trận trọng số:
\begin{equation}
\mathbf{W}_{masked} = \mathbf{W}_{df}\odot\mathcal{M}_1 ,
\label{eq:wmasked}
\end{equation}
sau đó mặt nạ nút được áp lên đầu ra, kèm các hệ số bù năng lượng:
\begin{equation}
\DF(\mathbf{Z};\mathbf{W}_{df},\mathbf{b}_{df}) = \frac{1}{1-p_2}\left[\frac{\mathbf{Z}\,\mathbf{W}_{masked}}{1-p_1}
+ \one_D\mathbf{b}_{df}^{\top}\right]\odot\mathcal{M}_2 ,
\label{eq:dropfusion}
\end{equation}
trong đó:
\begin{itemize}
  \item $\mathcal{M}_1\in\{0,1\}^{C_{in}\times C_{out}}$ là \emph{mặt nạ cạnh}. Mỗi phần tử bằng 0 với xác suất $p_1$,
        tức là cắt liên kết giữa một cặp kênh. Mặt nạ được lấy mẫu lại ở mỗi lô.
  \item $\mathcal{M}_2\in\{0,1\}^{D\times C_{out}}$ là \emph{mặt nạ nơ-ron}. Mỗi phần tử bằng 0 với xác suất $p_2$, tức
        là tắt đầu ra. Mặt nạ được lấy mẫu lại ở mỗi mẫu.
  \item Các hệ số $\frac{1}{1-p_1}$ và $\frac{1}{1-p_2}$ giữ cho kỳ vọng của tín hiệu không đổi giữa lúc huấn luyện và
        lúc kiểm thử.
\end{itemize}
Khi kiểm thử, $\mathcal{M}_1=\mathcal{M}_2=\one$, không nhân hệ số bù, và lớp trở thành
$\DF(\mathbf{Z})=\mathbf{Z}\mathbf{W}_{df}+\one_D\mathbf{b}_{df}^{\top}$.

CrossVBlock gồm hai lớp DropFusion theo cấu trúc nút thắt $C\to C'\to C$ (tham số
$\mathbf{W}_{df,1}\in\R^{C\times C'}$, $\mathbf{W}_{df,2}\in\R^{C'\times C}$), xen giữa là GELU:
\begin{equation}
\mathcal{X}_{cross} = \DF\Big(\GELU\big(\DF(\mathcal{X}'_{enc};\mathbf{W}_{df,1},\mathbf{b}_{df,1})\big);
\mathbf{W}_{df,2},\mathbf{b}_{df,2}\Big)\in\R^{D\times C}.
\label{eq:cross}
\end{equation}
Mỗi hàng của $\mathcal{X}'_{enc}$ là một đặc trưng (trong $D$ đặc trưng) của cả $C$ kênh. Phép nhân phải với
$\mathbf{W}_{df,1}$, $\mathbf{W}_{df,2}$ trộn thông tin \emph{giữa các kênh} trên từng đặc trưng đó.

\textbf{c) Khối Nội biến IntraVBlock.} Như DyTimeNet, một kết nối thặng dư cộng đặc trưng chéo với đặc trưng đã chuẩn
hóa, để thông tin gốc của từng kênh không bị mất:
\begin{equation}
\mathcal{X}_{intra} = \mathcal{X}_{cross} + \mathcal{X}'_{enc}\in\R^{D\times C}.
\label{eq:cross-res}
\end{equation}
Sau đó một mạng FFN hai lớp dùng chung được quét dọc theo từng kênh (từng cột), có LayerNorm phía trước và kết nối
thặng dư:
\begin{align}
\FFN(\mathbf{Z}) &= \Big(\operatorname{Drop}\big(\GELU(\mathbf{Z}^{\top}\mathbf{W}_1+\one_C\mathbf{b}_1^{\top})\big)\mathbf{W}_2
+ \one_C\mathbf{b}_2^{\top}\Big)^{\!\top}\in\R^{D\times C},
\label{eq:ffn}\\
\hat{\mathcal{X}}_{intra} &= \mathcal{X}_{intra} + \operatorname{Drop}\Big(\FFN\big(\LN(\mathcal{X}_{intra})\big)\Big)
\in\R^{D\times C},
\label{eq:intra}
\end{align}
trong đó $\mathbf{W}_1\in\R^{D\times D_{ff}}$, $\mathbf{W}_2\in\R^{D_{ff}\times D}$, và $\operatorname{Drop}$ là
Dropout thông thường với tỉ lệ $p_{drop}$. Phép chuyển vị $\mathbf{Z}^{\top}\in\R^{C\times D}$ đặt mỗi kênh thành một
hàng, nên $\mathbf{W}_1$, $\mathbf{W}_2$ chỉ biến đổi trên $D$ đặc trưng của từng kênh, không trộn các kênh.

\textbf{d) Chiếu ngược và Hiệu chỉnh cơ sở.} Biểu diễn ẩn được chiếu trở lại không gian cơ sở:
\begin{equation}
\Delta\bar{\mathbf{E}} = \mathbf{W}_{out}^{\top}\hat{\mathcal{X}}_{intra} + \mathbf{b}_{out}\one_C^{\top}\in\R^{PR\times C},
\qquad \mathbf{W}_{out}\in\R^{D\times PR}.
\label{eq:proj}
\end{equation}
Cột thứ $i$ được xếp lại thành $\Delta\mathbf{E}^{(i)}\in\R^{P\times R}$ (phép ngược của $\vecr$), rồi cộng vào cơ sở
ban đầu:
\begin{equation}
\mathbf{E}'^{(i)} = \mathbf{E}^{(i)} + \Delta\mathbf{E}^{(i)}.
\label{eq:basis-update}
\end{equation}
\textbf{Tác dụng và ý nghĩa:} $\mathbf{W}_{out}$ và $\mathbf{b}_{out}$ được khởi tạo bằng 0, nên ở thời điểm khởi tạo
$\Delta\mathbf{E}^{(i)}=\mathbf{0}$ và DyTimeBase trùng khớp với TimeBase. Trong quá trình huấn luyện, nhánh trộn kênh
chỉ thêm vào phần hiệu chỉnh nào giúp giảm sai số. Nhờ vậy mô hình không bắt đầu từ một mức tệ hơn xương sống độc
lập kênh.

\subsection{Khối Dự báo theo đoạn và Khôi phục thang đo}
\label{sec:forecast}

\textbf{a) Dự báo theo đoạn.} Một ánh xạ tuyến tính dùng chung sinh ra $N'$ đoạn tương lai từ $R$ chuỗi cơ sở đã
hiệu chỉnh:
\begin{equation}
\hat{\mathbf{Y}}^{(i)} = \mathbf{E}'^{(i)}\mathbf{W}_F + \one_P\mathbf{b}_F^{\top}\in\R^{P\times N'},
\qquad \mathbf{W}_F\in\R^{R\times N'},\;\mathbf{b}_F\in\R^{N'}.
\label{eq:forecast}
\end{equation}

\textbf{b) Khôi phục thang đo.} Trung bình theo pha ở~\eqref{eq:period-norm} được cộng trả lại cho mọi đoạn tương
lai:
\begin{equation}
\hat{\mathbf{Y}}'^{(i)} = \hat{\mathbf{Y}}^{(i)} + \bm{\mu}^{(i)}\one_{N'}^{\top}.
\label{eq:denorm}
\end{equation}

\textbf{c) Ghép đoạn.} Các cột được nối tiếp nhau thành chuỗi, và mô hình giữ lại $S$ bước đầu:
\begin{equation}
\hat{x}^{(i)}_{L+t} = \big[\hat{\mathbf{Y}}'^{(i)}\big]_{p,n}\quad\text{với } t=(n-1)P+p,\; t=1,\dots,S,
\qquad
\hat{\mathcal{X}} = \big[\hat{x}^{(i)}_{L+t}\big]_{t,i}\in\R^{S\times C}.
\label{eq:merge}
\end{equation}

\subsection{Huấn luyện mô hình}
\label{sec:training}

Tham số được chia thành hai nhóm:
\begin{itemize}
  \item Xương sống $\Theta_B=\{\mathbf{W}_B,\mathbf{b}_B,\mathbf{W}_F,\mathbf{b}_F\}$.
  \item Nhánh trộn kênh $\Theta_M$: các tham số của~\eqref{eq:embed}--\eqref{eq:proj}, gồm cả tham số của các lớp
        LayerNorm.
\end{itemize}
Cả hai nhóm được tối ưu đồng thời theo hàm mục tiêu~\eqref{eq:objective} bằng thuật toán Adam. Mỗi nhóm có một tốc
độ học riêng: $\alpha_B$ cho xương sống và $\alpha_M$ cho nhánh trộn kênh. Với $\mathbf{g}_k=\nabla_{\theta}\mathcal{L}$
tại bước $k$:
\begin{equation}
\mathbf{m}_k=\beta_1\mathbf{m}_{k-1}+(1-\beta_1)\mathbf{g}_k,\quad
\mathbf{v}_k=\beta_2\mathbf{v}_{k-1}+(1-\beta_2)\mathbf{g}_k^{2},\quad
\theta_{k+1}=\theta_k-\alpha\,\frac{\mathbf{m}_k/(1-\beta_1^{k})}{\sqrt{\mathbf{v}_k/(1-\beta_2^{k})}+\epsilon},
\label{eq:adam}
\end{equation}
trong đó $\alpha=\alpha_B$ nếu $\theta\in\Theta_B$ và $\alpha=\alpha_M$ nếu $\theta\in\Theta_M$.

Tốc độ học giữ nguyên trong 4 epoch đầu, sau đó giảm theo cấp số nhân với hệ số $0.8$ mỗi epoch. Lịch này áp dụng
cho cả hai nhóm:
\begin{equation}
\alpha^{(e)} = \begin{cases}\alpha, & e\le 4,\\ \alpha\cdot 0.8^{\,e-4}, & e>4,\end{cases}
\qquad e=1,2,\dots
\label{eq:lr-schedule}
\end{equation}
Huấn luyện chạy tối đa 30 epoch và dừng sớm nếu MSE trên tập Validation không giảm sau 5 epoch liên tiếp. Trọng số
có MSE Validation thấp nhất được dùng để đánh giá trên tập Test.

\subsection{Độ phức tạp tính toán}
\label{sec:complexity}

Số phép nhân--cộng (MACs) cho một mẫu là tổng của hai phần.

Xương sống, gồm~\eqref{eq:basis} và~\eqref{eq:forecast}:
\begin{equation}
\text{MACs}_{B} = C\,P\,R\,(N+N').
\label{eq:macs-backbone}
\end{equation}

Nhánh trộn kênh, gồm mã hóa và chiếu ngược, CrossVBlock và IntraVBlock:
\begin{equation}
\text{MACs}_{M} = \underbrace{2\,C\,PR\,D}_{\text{mã hóa, chiếu ngược}}
+ \underbrace{2\,D\,C\,C'}_{\text{CrossVBlock}}
+ \underbrace{2\,C\,D\,D_{ff}}_{\text{IntraVBlock}} .
\label{eq:macs-mix}
\end{equation}

Ví dụ với $C=321$, $L=S=720$, $P=24$ ($N=N'=30$), $R=6$, $D=64$, $C'=240$, $D_{ff}=128$:
\begin{itemize}
  \item Xương sống: $\text{MACs}_B\approx 2.8$M.
  \item Nhánh trộn kênh: $\text{MACs}_M\approx 5.9\text{M}+9.9\text{M}+5.3\text{M}$.
  \item Tổng cộng: khoảng $23.8$M MACs.
\end{itemize}
Để so sánh, DLinear cần $2CLS\approx 333$M MACs.

\subsection{Tóm tắt luồng dữ liệu}

\begin{table}[ht]
\centering
\caption{Kích thước dữ liệu qua từng khối của DyTimeBase (một mẫu)}
\label{tab:shapes}
\small
\begin{tabular}{@{}llll@{}}
\toprule
Khối & Phép biến đổi & Đầu ra & Công thức\\
\midrule
Input & -- & $\mathcal{X}_{train}\in\R^{L\times C}$ & \eqref{eq:input}\\
Phân đoạn & đệm, xếp $P\times N$ & $\mathbf{X}^{(i)}\in\R^{P\times N}$ & \eqref{eq:pad}, \eqref{eq:segment}\\
Chuẩn hóa theo pha & trừ $\bm{\mu}^{(i)}$ & $\tilde{\mathbf{X}}^{(i)}\in\R^{P\times N}$ & \eqref{eq:period-norm}\\
Trích xuất cơ sở & Linear $N\to R$ & $\mathbf{E}^{(i)}\in\R^{P\times R}$ & \eqref{eq:basis}, \eqref{eq:orth}\\
Mã hóa & Linear $PR\to D$ & $\mathcal{X}_{enc}\in\R^{D\times C}$ & \eqref{eq:embed}\\
CrossVBlock & LN, DF $C\to C'\to C$ & $\mathcal{X}_{cross}\in\R^{D\times C}$ & \eqref{eq:ln-cross}--\eqref{eq:cross}\\
IntraVBlock & thặng dư, LN, FFN $D\to 2D\to D$ & $\hat{\mathcal{X}}_{intra}\in\R^{D\times C}$ & \eqref{eq:cross-res}--\eqref{eq:intra}\\
Chiếu ngược & Linear $D\to PR$ (khởi tạo 0) & $\mathbf{E}'^{(i)}\in\R^{P\times R}$ & \eqref{eq:proj}, \eqref{eq:basis-update}\\
Dự báo theo đoạn & Linear $R\to N'$ & $\hat{\mathbf{Y}}^{(i)}\in\R^{P\times N'}$ & \eqref{eq:forecast}\\
Khôi phục, ghép & cộng $\bm{\mu}^{(i)}$, trải phẳng & $\hat{\mathcal{X}}\in\R^{S\times C}$ & \eqref{eq:denorm}, \eqref{eq:merge}\\
\bottomrule
\end{tabular}
\end{table}

\subsection{Các biến thể}

\begin{itemize}
  \item \textbf{Tham số riêng từng kênh} (\emph{individual}): thay $\mathbf{W}_B$, $\mathbf{b}_B$, $\mathbf{W}_F$, $\mathbf{b}_F$
        bằng $\mathbf{W}^{(i)}_B$, $\mathbf{b}^{(i)}_B$, $\mathbf{W}^{(i)}_F$, $\mathbf{b}^{(i)}_F$ cho từng kênh $i$. Số MACs
        không đổi, chỉ số tham số của xương sống tăng $C$ lần. Biến thể này dùng cho bộ dữ liệu Weather.
  \item \textbf{Chuẩn hóa toàn cửa sổ}: thay $\bm{\mu}^{(i)}$ trong~\eqref{eq:period-norm} bằng một trung bình vô hướng
        $\mu^{(i)}=\frac{1}{NP}\one_P^{\top}\mathbf{X}^{(i)}\one_N$.
  \item \textbf{Bỏ nhánh trộn kênh}: bỏ các bước từ~\eqref{eq:embed} đến~\eqref{eq:basis-update}, tức
        $\mathbf{E}'^{(i)}=\mathbf{E}^{(i)}$. Khi đó mô hình đúng bằng TimeBase.
\end{itemize}

%=====================================================================
\chapter{Thực nghiệm và kết quả}

Chương này đánh giá DyTimeBase trên 5 bộ dữ liệu chuẩn ở cả hai mặt: độ chính xác dự báo và chi phí tính toán. Các thực
nghiệm trả lời bốn câu hỏi:
\begin{enumerate}
  \item DyTimeBase chính xác hơn hay kém hơn các mô hình hiện đại ở những bộ dữ liệu nào?
  \item So với xương sống TimeBase trong cùng một quy trình huấn luyện, nhánh trộn kênh có cải thiện kết quả không?
  \item DyTimeBase nhanh và nhẹ đến mức nào so với các mô hình khác khi chạy trên cùng một thiết bị?
  \item Từng khối của nhánh trộn kênh (CrossVBlock, IntraVBlock) đóng góp bao nhiêu?
\end{enumerate}

%---------------------------------------------------------------------
\section{Mô tả dữ liệu thực nghiệm}

Thực nghiệm dùng lại 5 bộ dữ liệu của Đồ án I: ETTh1, ETTh2, ETTm1, ETTm2 (nhiệt độ dầu và phụ tải máy biến áp, 7 biến)
và Weather (21 chỉ số khí tượng). Đặc điểm chi tiết của từng bộ đã được mô tả ở Chương 3 của Đồ án I. Bảng~\ref{tab:data}
tóm tắt kích thước dữ liệu và số cửa sổ mẫu khi độ dài đầu vào là $L=720$.

\begin{table}[H]
\centering
\caption{Thống kê các bộ dữ liệu. Số cửa sổ tính với $L=720$, $S=96$ (huấn luyện / kiểm định / kiểm thử).}
\label{tab:data}
\small
\resizebox{\textwidth}{!}{%
\begin{tabular}{@{}lccccc@{}}
\toprule
Dataset & Số biến $C$ & Tần suất & Số bước dùng & Chia Train/Val/Test & Số cửa sổ\\
\midrule
ETTh1, ETTh2 & 7 & 1 giờ & 14\,400 & 12/4/4 tháng & 7\,825 / 2\,785 / 2\,785\\
ETTm1, ETTm2 & 7 & 15 phút & 57\,600 & 12/4/4 tháng & 33\,745 / 11\,425 / 11\,425\\
Weather & 21 & 10 phút & 52\,696 & 7 : 1 : 2 & 36\,072 / 5\,174 / 10\,445\\
\bottomrule
\end{tabular}}
\end{table}

%---------------------------------------------------------------------
\section{Thiết lập thực nghiệm}

\subsection{Tiền xử lý dữ liệu}

Quy trình tiền xử lý giống Đồ án I, chỉ khác độ dài đầu vào:
\begin{itemize}
  \item \textbf{Chuẩn hóa Z-score} theo từng biến, với trung bình và độ lệch chuẩn chỉ tính trên tập huấn luyện rồi dùng
        lại cho tập kiểm định và kiểm thử.
  \item \textbf{Cửa sổ trượt}: mỗi mẫu gồm $L=720$ bước quá khứ và $S\in\{96,192,336,720\}$ bước tương lai. Độ dài
        $L=720$ được chọn để trùng với thiết lập của TimeBase~\cite{timebase}, nhờ đó số liệu của các mô hình so sánh
        lấy từ~\cite{timebase} có cùng điều kiện đầu vào.
  \item \textbf{Chia dữ liệu theo thời gian}: tập kiểm định và kiểm thử được lấy thêm $L$ bước phía trước làm đầu vào,
        nhưng toàn bộ nhãn dự báo nằm trong đúng tập tương ứng. Khi đánh giá, mọi cửa sổ của tập kiểm thử đều được dùng
        (không bỏ lô cuối).
\end{itemize}

\subsection{Siêu tham số và quá trình huấn luyện}

Siêu tham số của DyTimeBase gồm hai nhóm (Bảng~\ref{tab:hparams}):
\begin{itemize}
  \item \textbf{Xương sống} ($P$, $R$, tốc độ học $\alpha_B$, kích thước lô, $\lambda_{orth}$, tham số riêng từng kênh):
        lấy nguyên theo script chính thức của TimeBase cho từng bộ dữ liệu và từng độ dài dự báo.
  \item \textbf{Nhánh trộn kênh} ($D$, $p_1$, $p_2$): giữ nguyên bộ siêu tham số đã tối ưu bằng Grid Search cho DyTimeNet
        trong Đồ án I (Bảng 3.1). DyTimeBase \emph{không} được tinh chỉnh thêm.
\end{itemize}

\begin{table}[H]
\centering
\caption{Siêu tham số của DyTimeBase. Xương sống: theo script chính thức của TimeBase~\cite{timebase} cho từng độ dài dự báo $S\in\{96,192,336,720\}$. Nhánh trộn kênh: giữ nguyên bộ đã tối ưu của DyTimeNet (Đồ án I, Bảng 3.1).}
\label{tab:hparams}
\small
\resizebox{\textwidth}{!}{\begin{tabular}{@{}llllllccc@{}}
\toprule
 & \multicolumn{5}{c}{Xương sống (theo $S=96/192/336/720$)} & & \multicolumn{2}{c}{Nhánh trộn kênh}\\
\cmidrule(lr){2-6}\cmidrule(l){8-9}
Dataset & $P$ & $R$ & $\alpha_B$ & Batch & $\lambda_{orth}$ & Riêng kênh & $D$ & $p_1$ / $p_2$\\
\midrule
ETTh1 & 24/24/24/24 & 6/6/6/6 & 0.1/0.4/0.4/0.05 & 64/256/256/64 & 0.16/0.16/0.08/0.12 & không & 64 & 0.1 / 0.1\\
ETTh2 & 24/24/24/24 & 6/4/6/6 & 0.2/0.06/0.4/0.4 & 512/512/64/64 & 0.2/0.08/0.12/0.12 & không & 64 & 0.6 / 0.6\\
ETTm1 & 4/4/4/6 & 18/20/20/20 & 0.02/0.02/0.02/0.01 & 512/256/256/128 & 0.04/0.04/0.08/0.12 & không & 64 & 0.1 / 0.2\\
ETTm2 & 4/4/4/6 & 20/20/20/20 & 0.01/0.01/0.01/0.01 & 64/64/64/64 & 0.04/0.04/0.04/0.04 & không & 32 & 0.1 / 0.1\\
Weather & 4/4/4/4 & 6/6/6/6 & 0.02/0.02/0.05/0.05 & 512/512/512/512 & 0.04/0.04/0.08/0.04 & có & 64 & 0.1 / 0.1\\
\bottomrule
\end{tabular}}
\end{table}


Các thiết lập chung cho mọi bộ dữ liệu:
\begin{itemize}
  \item Hàm mục tiêu $\mathcal{L}=\mathcal{L}_{MSE}+\lambda_{orth}\mathcal{L}_{orth}$~\eqref{eq:objective}, tối ưu bằng Adam
        không dùng weight decay và không cắt gradient.
  \item Tốc độ học của nhánh trộn kênh $\alpha_M=5\cdot10^{-4}$ (bằng tốc độ học của DyTimeNet trong Đồ án I); lịch tốc
        độ học theo~\eqref{eq:lr-schedule} áp dụng cho cả hai nhóm tham số.
  \item $C'=\lfloor 0.75\,C\rfloor$, $D_{ff}=2D$, Dropout trong IntraVBlock $p_{drop}=0.1$.
  \item Tối đa 30 epoch, dừng sớm khi MSE trên tập kiểm định không giảm sau 5 epoch; trọng số có MSE kiểm định thấp nhất
        được dùng để đánh giá trên tập kiểm thử.
  \item Mỗi cấu hình chạy với 3 seed (2021, 2022, 2023); kết quả báo cáo là trung bình của 3 lần chạy.
\end{itemize}

\subsection{Các mô hình so sánh}

DyTimeBase được so sánh với 14 mô hình, chia làm ba nhóm:
\begin{itemize}
  \item \textbf{TimeBase$^{*}$}: TimeBase~\cite{timebase} cài đặt lại theo mã nguồn chính thức và chạy trong \emph{cùng quy
        trình} với DyTimeBase (cùng dữ liệu, cùng siêu tham số xương sống, cùng 3 seed). Với cùng trọng số, đầu ra của bản
        cài đặt lại lệch so với mã nguồn chính thức không quá $3\cdot10^{-7}$. Đây là phép so sánh công bằng nhất, vì hai mô
        hình chỉ khác nhau đúng ở nhánh trộn kênh.
  \item \textbf{CycleNet và TQNet}: CycleNet/Linear, CycleNet/MLP~\cite{cyclenet} (học chu kỳ tuần hoàn, độc lập kênh) và
        TQNet~\cite{tqnet} (truy vấn theo chu kỳ, có học quan hệ giữa các kênh) là hai mô hình nhẹ gần đây có code công
        khai. Cả hai được cài đặt lại theo mã nguồn chính thức (với cùng trọng số, đầu ra trùng khớp hoàn toàn) và chạy
        trong cùng quy trình, 3 seed, với siêu tham số chính thức của từng mô hình: CycleNet theo script $L=720$ của tác
        giả; TQNet chỉ có script $L=96$ nên giữ nguyên siêu tham số và chỉ đổi độ dài đầu vào thành 720.
  \item \textbf{Số liệu công bố}: TimeBase$^{\dagger}$, SparseTSF~\cite{sparsetsf}, TimeMixer~\cite{timemixer},
        FITS~\cite{fits}, iTransformer~\cite{itransformer}, DLinear~\cite{dlinear}, PatchTST~\cite{patchtst},
        TimesNet~\cite{timesnet}, FEDformer~\cite{fedformer}, Autoformer~\cite{autoformer} và Informer~\cite{informer}. Kết
        quả lấy từ Table 10 của~\cite{timebase}, nơi mọi mô hình được chạy với cùng $L=720$ và cùng cách chia dữ liệu.
        Các mô hình Transformer cỡ lớn không được huấn luyện lại vì chi phí trên CPU quá cao (ví dụ, một epoch của
        PatchTST ở cấu hình Electricity mất khoảng 21 giờ, Bảng~\ref{tab:efficiency}).
\end{itemize}

\subsection{Môi trường thực nghiệm}

Toàn bộ huấn luyện và đo đạc chạy trên một máy tính xách tay không có GPU: CPU AMD Ryzen 7 7840HS (8 nhân), 16 GB RAM (15.2 GB khả dụng),
Windows 11, Python 3.13, PyTorch 2.6.0 (bản CPU).

%---------------------------------------------------------------------
\section{Các chỉ số đánh giá}

\subsection{Chỉ số độ chính xác}

Độ chính xác được đo bằng MSE và MAE như các công thức (3.1), (3.2) của Đồ án I:
\begin{equation}
\text{MSE}=\frac{1}{n}\sum_{i=1}^{n}(y_i-\hat{y}_i)^2,\qquad
\text{MAE}=\frac{1}{n}\sum_{i=1}^{n}\lvert y_i-\hat{y}_i\rvert ,
\label{eq:metrics}
\end{equation}
trong đó tổng lấy trên mọi cửa sổ kiểm thử, mọi bước dự báo và mọi biến. Sai số được tính trên thang đã chuẩn hóa
Z-score, theo quy ước của các nghiên cứu được so sánh.

\subsection{Chỉ số hiệu năng}

Theo~\cite{timebase}, hiệu năng được đo bằng năm chỉ số trên cấu hình Electricity ($C=321$, $L=S=720$):
\begin{itemize}
  \item \textbf{Params}: số tham số huấn luyện được.
  \item \textbf{MACs}: số phép nhân--cộng để dự báo một mẫu.
  \item \textbf{Max Mem}: bộ nhớ đỉnh khi huấn luyện với lô 12 mẫu. Vì thiết bị không có GPU, đây là bộ nhớ RAM của
        tiến trình, không phải bộ nhớ GPU như trong~\cite{timebase}.
  \item \textbf{Epoch Time}: thời gian một epoch huấn luyện, tính bằng thời gian một bước (đo thực tế) nhân với số bước của
        một epoch.
  \item \textbf{Infer Time}: thời gian dự báo một mẫu trên CPU, dùng 1 luồng. Các mô hình nhẹ được đo 5 lần xen kẽ và lấy
        trung vị.
\end{itemize}
Mỗi mô hình được đo trong một tiến trình riêng, với dữ liệu ngẫu nhiên cùng kích thước (tốc độ không phụ thuộc giá
trị dữ liệu). Các mô hình so sánh được dựng từ mã nguồn chính thức: Time-Series-Library~\cite{tslib} với cấu hình script
chính thức, PatchTST theo kho mã của~\cite{patchtst}, FITS theo kho mã của~\cite{fits}.

%---------------------------------------------------------------------
\section{Kết quả thực nghiệm}

\subsection{So sánh độ chính xác}

Bảng~\ref{tab:acc-mse} và Bảng~\ref{tab:acc-mae} trình bày MSE và MAE của 16 mô hình trên 5 bộ dữ liệu.

\begin{sidewaystable}
\centering
\caption{Kết quả dự báo dài hạn (MSE), độ dài đầu vào $L=720$ cho mọi mô hình. DyTimeBase và TimeBase$^{*}$: tự chạy, trung bình 3 seed. Các cột còn lại (TimeBase$^{\dagger}$ là số công bố): trích từ Table 10 của~\cite{timebase}. \textbf{Đậm}: tốt nhất; \underline{gạch chân}: thứ hai.}
\label{tab:acc-mse}
\scriptsize
\setlength{\tabcolsep}{3.2pt}
\begin{tabular}{@{}llccccccccccccc@{}}
\toprule
Dataset & $S$ & \textbf{DyTimeBase} & TimeBase$^{*}$ & TimeBase$^{\dagger}$ & SparseTSF & TimeMixer & FITS & iTransformer & DLinear & PatchTST & TimesNet & FEDformer & Autoformer & Informer\\
\midrule
\multirow{5}{*}{ETTh1} & 96 & \textbf{0.347} & 0.358 & \underline{0.349} & 0.362 & 0.410 & 0.380 & 0.389 & 0.378 & 0.377 & 0.437 & 0.485 & 0.555 & 1.269\\
 & 192 & \underline{0.388} & 0.396 & \textbf{0.387} & 0.404 & 0.448 & 0.415 & 0.424 & 0.415 & 0.413 & 0.456 & 0.481 & 0.599 & 1.487\\
 & 336 & \underline{0.420} & 0.423 & \textbf{0.408} & 0.435 & 0.482 & 0.439 & 0.456 & 0.449 & 0.436 & 0.494 & 0.522 & 0.853 & 1.544\\
 & 720 & 0.456 & 0.466 & 0.439 & \textbf{0.426} & 0.475 & \underline{0.433} & 0.545 & 0.507 & 0.455 & 0.632 & 0.604 & 0.899 & 1.481\\
 & TB & \underline{0.403} & 0.411 & \textbf{0.396} & 0.407 & 0.454 & 0.417 & 0.454 & 0.437 & 0.420 & 0.505 & 0.523 & 0.727 & 1.445\\
\midrule
\multirow{5}{*}{ETTh2} & 96 & 0.286 & 0.295 & 0.292 & 0.294 & 0.315 & \textbf{0.271} & 0.305 & 0.294 & \underline{0.276} & 0.349 & 0.401 & 0.541 & 5.189\\
 & 192 & 0.347 & 0.352 & \underline{0.339} & 0.340 & 0.383 & \textbf{0.332} & 0.405 & 0.412 & 0.342 & 0.500 & 0.425 & 1.207 & 6.514\\
 & 336 & 0.378 & 0.390 & \underline{0.358} & 0.360 & 0.385 & \textbf{0.355} & 0.411 & 0.471 & 0.364 & 0.445 & 0.427 & 0.825 & 5.284\\
 & 720 & 0.421 & 0.413 & 0.400 & \underline{0.383} & 0.432 & \textbf{0.378} & 0.448 & 0.740 & 0.395 & 0.438 & 0.462 & 1.772 & 4.955\\
 & TB & 0.358 & 0.363 & 0.347 & \underline{0.344} & 0.379 & \textbf{0.334} & 0.392 & 0.479 & \underline{0.344} & 0.433 & 0.429 & 1.086 & 5.486\\
\midrule
\multirow{5}{*}{ETTm1} & 96 & \textbf{0.296} & 0.318 & 0.311 & 0.314 & 0.332 & 0.313 & 0.315 & 0.314 & \underline{0.298} & 0.359 & 0.406 & 0.455 & 0.632\\
 & 192 & \textbf{0.328} & 0.363 & 0.338 & 0.348 & 0.355 & 0.339 & 0.349 & 0.347 & \underline{0.335} & 0.368 & 0.450 & 0.562 & 1.131\\
 & 336 & \textbf{0.364} & 0.372 & \textbf{0.364} & 0.368 & 0.386 & 0.367 & 0.381 & 0.367 & \underline{0.366} & 0.429 & 0.436 & 0.737 & 1.391\\
 & 720 & 0.416 & 0.418 & \textbf{0.413} & 0.419 & 0.452 & 0.417 & 0.437 & \underline{0.415} & 0.420 & 0.477 & 0.462 & 0.503 & 1.397\\
 & TB & \textbf{0.351} & 0.367 & 0.356 & 0.362 & 0.381 & 0.359 & 0.370 & 0.361 & \underline{0.355} & 0.408 & 0.439 & 0.564 & 1.138\\
\midrule
\multirow{5}{*}{ETTm2} & 96 & 0.167 & 0.166 & \textbf{0.162} & 0.167 & 0.192 & 0.166 & 0.179 & \underline{0.163} & 0.165 & 0.200 & 0.339 & 0.325 & 1.870\\
 & 192 & 0.237 & \textbf{0.218} & \textbf{0.218} & \underline{0.219} & 0.253 & \textbf{0.218} & 0.239 & 0.223 & \underline{0.219} & 0.274 & 0.397 & 0.369 & 2.807\\
 & 336 & 0.284 & 0.277 & \underline{0.270} & 0.271 & 0.307 & 0.271 & 0.309 & 0.291 & \textbf{0.268} & 0.340 & 0.449 & 0.418 & 4.442\\
 & 720 & 0.366 & 0.355 & \textbf{0.352} & \underline{0.353} & 0.380 & \textbf{0.352} & 0.387 & 0.407 & \textbf{0.352} & 0.384 & 0.451 & 0.612 & 5.258\\
 & TB & 0.263 & 0.254 & \textbf{0.251} & 0.253 & 0.283 & \underline{0.252} & 0.278 & 0.271 & \textbf{0.251} & 0.299 & 0.409 & 0.431 & 3.594\\
\midrule
\multirow{5}{*}{Weather} & 96 & \textbf{0.143} & \underline{0.144} & 0.146 & 0.174 & 0.163 & 0.176 & 0.159 & 0.174 & 0.149 & 0.176 & 0.289 & 0.323 & 0.283\\
 & 192 & \underline{0.187} & \underline{0.187} & \textbf{0.185} & 0.216 & 0.201 & 0.216 & 0.203 & 0.215 & 0.193 & 0.219 & 0.340 & 0.389 & 0.445\\
 & 336 & 0.240 & \underline{0.238} & \textbf{0.236} & 0.260 & 0.258 & 0.261 & 0.253 & 0.262 & 0.240 & 0.277 & 0.370 & 0.497 & 0.587\\
 & 720 & 0.314 & \textbf{0.309} & \textbf{0.309} & 0.325 & 0.329 & 0.325 & 0.317 & 0.319 & \underline{0.312} & 0.344 & 0.420 & 0.573 & 0.953\\
 & TB & \underline{0.221} & \textbf{0.219} & \textbf{0.219} & 0.244 & 0.238 & 0.244 & 0.233 & 0.242 & 0.224 & 0.254 & 0.355 & 0.446 & 0.567\\
\bottomrule
\end{tabular}
\end{sidewaystable}

\begin{sidewaystable}
\centering
\caption{Kết quả dự báo dài hạn (MAE), độ dài đầu vào $L=720$ cho mọi mô hình. DyTimeBase và TimeBase$^{*}$: tự chạy, trung bình 3 seed. Các cột còn lại (TimeBase$^{\dagger}$ là số công bố): trích từ Table 10 của~\cite{timebase}. \textbf{Đậm}: tốt nhất; \underline{gạch chân}: thứ hai.}
\label{tab:acc-mae}
\scriptsize
\setlength{\tabcolsep}{3.2pt}
\begin{tabular}{@{}llccccccccccccc@{}}
\toprule
Dataset & $S$ & \textbf{DyTimeBase} & TimeBase$^{*}$ & TimeBase$^{\dagger}$ & SparseTSF & TimeMixer & FITS & iTransformer & DLinear & PatchTST & TimesNet & FEDformer & Autoformer & Informer\\
\midrule
\multirow{5}{*}{ETTh1} & 96 & \underline{0.386} & \textbf{0.384} & \textbf{0.384} & 0.389 & 0.441 & 0.402 & 0.421 & 0.402 & 0.408 & 0.454 & 0.500 & 0.558 & 0.855\\
 & 192 & \textbf{0.408} & \underline{0.410} & \underline{0.410} & 0.412 & 0.465 & 0.424 & 0.446 & 0.425 & 0.431 & 0.469 & 0.498 & 0.575 & 0.943\\
 & 336 & 0.428 & \underline{0.421} & \textbf{0.418} & 0.428 & 0.490 & 0.439 & 0.469 & 0.449 & 0.446 & 0.494 & 0.521 & 0.702 & 0.945\\
 & 720 & 0.461 & 0.466 & \textbf{0.446} & \underline{0.448} & 0.500 & 0.457 & 0.532 & 0.517 & 0.475 & 0.578 & 0.575 & 0.730 & 0.975\\
 & TB & 0.421 & 0.420 & \textbf{0.414} & \underline{0.419} & 0.474 & 0.430 & 0.467 & 0.448 & 0.440 & 0.499 & 0.523 & 0.641 & 0.929\\
\midrule
\multirow{5}{*}{ETTh2} & 96 & 0.344 & 0.354 & 0.345 & 0.346 & 0.380 & \textbf{0.336} & 0.361 & 0.360 & \underline{0.339} & 0.403 & 0.451 & 0.559 & 1.812\\
 & 192 & 0.388 & 0.393 & 0.387 & \underline{0.377} & 0.415 & \textbf{0.374} & 0.421 & 0.437 & 0.385 & 0.488 & 0.464 & 0.866 & 2.011\\
 & 336 & 0.431 & 0.435 & 0.410 & \underline{0.398} & 0.438 & \textbf{0.396} & 0.436 & 0.478 & 0.405 & 0.465 & 0.471 & 0.719 & 1.859\\
 & 720 & 0.462 & 0.460 & 0.448 & \underline{0.425} & 0.471 & \textbf{0.423} & 0.470 & 0.609 & 0.434 & 0.465 & 0.493 & 1.062 & 1.884\\
 & TB & 0.406 & 0.410 & 0.397 & \underline{0.387} & 0.426 & \textbf{0.382} & 0.422 & 0.471 & 0.391 & 0.455 & 0.470 & 0.801 & 1.891\\
\midrule
\multirow{5}{*}{ETTm1} & 96 & 0.352 & 0.361 & \underline{0.351} & 0.359 & 0.384 & 0.357 & 0.369 & \textbf{0.350} & 0.352 & 0.391 & 0.441 & 0.464 & 0.574\\
 & 192 & \underline{0.371} & 0.393 & \underline{0.371} & 0.376 & 0.398 & \textbf{0.369} & 0.388 & 0.381 & 0.373 & 0.398 & 0.477 & 0.514 & 0.802\\
 & 336 & 0.394 & 0.394 & \underline{0.386} & \underline{0.386} & 0.416 & \textbf{0.385} & 0.409 & 0.387 & 0.394 & 0.438 & 0.466 & 0.608 & 0.923\\
 & 720 & 0.428 & 0.418 & \underline{0.414} & \textbf{0.413} & 0.457 & 0.417 & 0.439 & 0.415 & 0.421 & 0.474 & 0.479 & 0.502 & 0.973\\
 & TB & 0.386 & 0.391 & \textbf{0.381} & 0.384 & 0.414 & \underline{0.382} & 0.401 & 0.383 & 0.385 & 0.425 & 0.466 & 0.522 & 0.818\\
\midrule
\multirow{5}{*}{ETTm2} & 96 & 0.262 & \underline{0.257} & \textbf{0.256} & 0.259 & 0.285 & \textbf{0.256} & 0.274 & \underline{0.257} & 0.260 & 0.288 & 0.406 & 0.391 & 1.002\\
 & 192 & 0.314 & \textbf{0.293} & \textbf{0.293} & \underline{0.297} & 0.329 & \textbf{0.293} & 0.314 & 0.304 & 0.298 & 0.337 & 0.452 & 0.414 & 1.314\\
 & 336 & 0.346 & 0.334 & \textbf{0.328} & \underline{0.330} & 0.362 & \textbf{0.328} & 0.356 & 0.355 & 0.333 & 0.382 & 0.491 & 0.452 & 1.661\\
 & 720 & 0.396 & \underline{0.383} & \textbf{0.380} & \textbf{0.380} & 0.412 & \textbf{0.380} & 0.407 & 0.433 & 0.386 & 0.407 & 0.499 & 0.594 & 1.914\\
 & TB & 0.330 & \underline{0.317} & \textbf{0.314} & \underline{0.317} & 0.347 & \textbf{0.314} & 0.338 & 0.337 & 0.319 & 0.353 & 0.462 & 0.463 & 1.473\\
\midrule
\multirow{5}{*}{Weather} & 96 & 0.205 & \underline{0.199} & \textbf{0.198} & 0.231 & 0.223 & 0.232 & 0.212 & 0.242 & \underline{0.199} & 0.234 & 0.342 & 0.389 & 0.361\\
 & 192 & 0.250 & \underline{0.242} & \textbf{0.241} & 0.267 & 0.254 & 0.268 & 0.252 & 0.277 & 0.243 & 0.270 & 0.394 & 0.423 & 0.461\\
 & 336 & 0.300 & \underline{0.285} & \textbf{0.281} & 0.299 & 0.300 & 0.299 & 0.291 & 0.319 & \textbf{0.281} & 0.311 & 0.408 & 0.495 & 0.526\\
 & 720 & 0.351 & 0.335 & \textbf{0.331} & 0.345 & 0.348 & 0.346 & 0.337 & 0.359 & \underline{0.334} & 0.356 & 0.421 & 0.520 & 0.703\\
 & TB & 0.277 & 0.265 & \textbf{0.263} & 0.285 & 0.281 & 0.286 & 0.273 & 0.299 & \underline{0.264} & 0.293 & 0.391 & 0.457 & 0.513\\
\bottomrule
\end{tabular}
\end{sidewaystable}


Xét MSE trung bình trên 4 độ dài dự báo:
\begin{itemize}
  \item Trên \textbf{ETTm1}, DyTimeBase đạt MSE thấp nhất trong cả 16 mô hình (0.351), tốt hơn PatchTST (0.355), số
        công bố của TimeBase (0.356) và CycleNet/Linear (0.358).
  \item Trên \textbf{ETTh1} và \textbf{Weather}, DyTimeBase đứng thứ hai (0.403 và 0.221), chỉ sau TimeBase (0.396 và
        0.219), và tốt hơn PatchTST, SparseTSF, FITS, CycleNet, TQNet cùng mọi mô hình Transformer.
  \item Trên \textbf{ETTh2} và \textbf{ETTm2}, DyTimeBase (0.358 và 0.263) kém nhóm mô hình tuyến tính mạnh (FITS,
        SparseTSF, TimeBase, CycleNet/Linear) và PatchTST. CycleNet/Linear đạt MSE thấp nhất trên ETTm2 (0.249).
\end{itemize}

Bảng~\ref{tab:h2h} đếm số lần DyTimeBase thắng từng mô hình trên 20 cặp (bộ dữ liệu, độ dài dự báo). DyTimeBase thắng
toàn bộ 20/20 trước TimeMixer, iTransformer, TimesNet, FEDformer, Autoformer và Informer, và thắng DLinear 17/20. So với
TQNet, mô hình nhẹ duy nhất trong bảng cũng học quan hệ giữa các kênh, DyTimeBase thắng 16/20 và có MSE trung bình thấp
hơn trên cả 5 bộ dữ liệu; tỉ số với CycleNet/MLP cũng là 16/20. Với nhóm mô hình mạnh nhất (SparseTSF, FITS, PatchTST,
CycleNet/Linear), DyTimeBase ở mức ngang ngửa: thắng trên ETTh1, ETTm1, Weather và thua trên ETTh2, ETTm2. DyTimeBase vẫn
kém số liệu công bố của TimeBase (thắng 5/20).

\begin{table}[H]
\centering
\caption{DyTimeBase đối đầu từng mô hình theo MSE: số ô thắng / hòa / thua trên 20 ô (5 dataset $\times$ 4 độ dài dự báo, làm tròn 3 chữ số) và các dataset DyTimeBase có MSE trung bình tốt hơn hoặc kém hơn.}
\label{tab:h2h}
\small
\resizebox{\textwidth}{!}{\begin{tabular}{@{}lcccll@{}}
\toprule
Mô hình & Thắng & Hòa & Thua & DyTimeBase tốt hơn & DyTimeBase kém hơn\\
\midrule
TimeBase$^{*}$ & 12 & 1 & 7 & ETTh1, ETTh2, ETTm1 & ETTm2, Weather\\
TimeBase$^{\dagger}$ & 5 & 1 & 14 & ETTm1 & ETTh1, ETTh2, ETTm2, Weather\\
SparseTSF & 12 & 1 & 7 & ETTh1, ETTm1, Weather & ETTh2, ETTm2\\
TimeMixer & 20 & 0 & 0 & ETTh1, ETTh2, ETTm1, ETTm2, Weather & --\\
FITS & 11 & 0 & 9 & ETTh1, ETTm1, Weather & ETTh2, ETTm2\\
iTransformer & 20 & 0 & 0 & ETTh1, ETTh2, ETTm1, ETTm2, Weather & --\\
DLinear & 17 & 0 & 3 & ETTh1, ETTh2, ETTm1, ETTm2, Weather & --\\
PatchTST & 9 & 1 & 10 & ETTh1, ETTm1, Weather & ETTh2, ETTm2\\
TimesNet & 20 & 0 & 0 & ETTh1, ETTh2, ETTm1, ETTm2, Weather & --\\
FEDformer & 20 & 0 & 0 & ETTh1, ETTh2, ETTm1, ETTm2, Weather & --\\
Autoformer & 20 & 0 & 0 & ETTh1, ETTh2, ETTm1, ETTm2, Weather & --\\
Informer & 20 & 0 & 0 & ETTh1, ETTh2, ETTm1, ETTm2, Weather & --\\
\bottomrule
\end{tabular}}
\end{table}


\subsection{So sánh với TimeBase trong cùng quy trình}

Số liệu công bố của TimeBase được chạy một lần, trong khi DyTimeBase là trung bình 3 seed. Bản TimeBase$^{*}$ chạy lại
trong cùng quy trình cho kết quả kém hơn số công bố trên ETTh1, ETTh2 và ETTm1 (ví dụ 0.411 so với 0.396 trên ETTh1),
nên so sánh DyTimeBase với số công bố là bất lợi cho DyTimeBase. Bảng~\ref{tab:vs-timebase} so sánh trực tiếp với
TimeBase$^{*}$, nơi hai mô hình chỉ khác nhau ở nhánh trộn kênh.

\begin{table}[H]
\centering
\caption{DyTimeBase so với TimeBase$^{*}$ chạy lại trong cùng quy trình (cùng xương sống, cùng siêu tham số xương sống). MSE trung bình $\pm$ độ lệch chuẩn qua 3 seed; cột cuối là chênh lệch tương đối.}
\label{tab:vs-timebase}
\small
\begin{tabular}{@{}llccr@{}}
\toprule
Dataset & $S$ & DyTimeBase & TimeBase$^{*}$ & Chênh lệch\\
\midrule
\multirow{5}{*}{ETTh1} & 96 & \textbf{0.3474 $\pm$ 0.0014} & 0.3575 $\pm$ 0.0013 & $-2.8\%$\\
 & 192 & \textbf{0.3882 $\pm$ 0.0014} & 0.3963 $\pm$ 0.0026 & $-2.0\%$\\
 & 336 & \textbf{0.4204 $\pm$ 0.0051} & 0.4226 $\pm$ 0.0006 & $-0.5\%$\\
 & 720 & \textbf{0.4558 $\pm$ 0.0213} & 0.4656 $\pm$ 0.0080 & $-2.1\%$\\
 & TB & \textbf{0.4030} & 0.4105 & $-1.8\%$\\
\midrule
\multirow{5}{*}{ETTh2} & 96 & \textbf{0.2856 $\pm$ 0.0021} & 0.2953 $\pm$ 0.0004 & $-3.3\%$\\
 & 192 & \textbf{0.3467 $\pm$ 0.0029} & 0.3525 $\pm$ 0.0022 & $-1.7\%$\\
 & 336 & \textbf{0.3784 $\pm$ 0.0027} & 0.3898 $\pm$ 0.0094 & $-2.9\%$\\
 & 720 & 0.4212 $\pm$ 0.0046 & \textbf{0.4134 $\pm$ 0.0074} & $+1.9\%$\\
 & TB & \textbf{0.3580} & 0.3627 & $-1.3\%$\\
\midrule
\multirow{5}{*}{ETTm1} & 96 & \textbf{0.2961 $\pm$ 0.0031} & 0.3177 $\pm$ 0.0019 & $-6.8\%$\\
 & 192 & \textbf{0.3276 $\pm$ 0.0003} & 0.3626 $\pm$ 0.0219 & $-9.6\%$\\
 & 336 & \textbf{0.3642 $\pm$ 0.0065} & 0.3717 $\pm$ 0.0042 & $-2.0\%$\\
 & 720 & \textbf{0.4156 $\pm$ 0.0060} & 0.4177 $\pm$ 0.0027 & $-0.5\%$\\
 & TB & \textbf{0.3509} & 0.3674 & $-4.5\%$\\
\midrule
\multirow{5}{*}{ETTm2} & 96 & 0.1665 $\pm$ 0.0083 & \textbf{0.1660 $\pm$ 0.0001} & $+0.3\%$\\
 & 192 & 0.2369 $\pm$ 0.0043 & \textbf{0.2182 $\pm$ 0.0001} & $+8.6\%$\\
 & 336 & 0.2837 $\pm$ 0.0086 & \textbf{0.2767 $\pm$ 0.0106} & $+2.5\%$\\
 & 720 & 0.3663 $\pm$ 0.0038 & \textbf{0.3549 $\pm$ 0.0033} & $+3.2\%$\\
 & TB & 0.2634 & \textbf{0.2539} & $+3.7\%$\\
\midrule
\multirow{5}{*}{Weather} & 96 & \textbf{0.1435 $\pm$ 0.0016} & 0.1441 $\pm$ 0.0003 & $-0.4\%$\\
 & 192 & 0.1867 $\pm$ 0.0009 & \textbf{0.1865 $\pm$ 0.0000} & $+0.1\%$\\
 & 336 & 0.2403 $\pm$ 0.0009 & \textbf{0.2382 $\pm$ 0.0001} & $+0.9\%$\\
 & 720 & 0.3141 $\pm$ 0.0063 & \textbf{0.3092 $\pm$ 0.0003} & $+1.6\%$\\
 & TB & 0.2211 & \textbf{0.2195} & $+0.8\%$\\
\bottomrule
\end{tabular}
\end{table}


Nhánh trộn kênh giảm MSE trung bình trên ETTh1 ($-1.8\%$), ETTh2 ($-1.3\%$) và ETTm1 ($-4.5\%$), nhưng tăng MSE trên
ETTm2 ($+3.7\%$) và Weather ($+0.8\%$). Mức cải thiện lớn nhất xuất hiện ở các độ dài dự báo ngắn (ETTm1 với $S=96$ và
$S=192$: $-6.8\%$ và $-9.6\%$). Trên ETTm2, các lần chạy của DyTimeBase thường dừng sớm sau 6--7 epoch, cho thấy nhánh
phi tuyến dễ quá khớp trên chuỗi tương đối trơn, nơi một mô hình tuyến tính đã đủ tốt.

\subsection{Hiệu năng tính toán}

\begin{table}[H]
\centering
\caption{Hiệu năng tính toán trên cùng một máy CPU, cấu hình như Table~3 của~\cite{timebase}: Electricity ($C=321$), $L=S=720$. Max Mem: đỉnh bộ nhớ RAM khi huấn luyện với batch 12; Epoch Time: thời gian một bước huấn luyện (batch 12) nhân số bước của một epoch (1\,414); Infer Time: thời gian dự báo một mẫu trên CPU (1 luồng), như~\cite{timebase}; các mô hình nhẹ lấy trung vị của 5 lần đo. Sắp xếp theo Infer Time.}
\label{tab:efficiency}
\small
\resizebox{\textwidth}{!}{\begin{tabular}{@{}lrrrrrr@{}}
\toprule
Mô hình & Params & MACs & Max Mem (MB) & Epoch Time (s) & Infer Time (ms) & So với DyTimeBase\\
\midrule
TimeBase & 0.40K & 2.77M & 117 & 24 & 1.01 & 0.62$\times$\\
\textbf{DyTimeBase} & 190.51K & 23.81M & 146 & 40 & 1.63 & 1.00$\times$\\
SparseTSF & 0.93K & 12.71M & 275 & 54 & 2.39 & 1.46$\times$\\
CycleNet-L & 573.05K & 166.41M & 153 & 74 & 6.71 & 4.11$\times$\\
FITS & 133.64K & 42.73M & 210 & 81 & 7.48 & 4.58$\times$\\
DLinear & 1.04M & 332.81M & 150 & 58 & 7.88 & 4.83$\times$\\
CycleNet-M & 792.44K & 236.67M & 162 & 92 & 7.93 & 4.86$\times$\\
DyTimeNet (Đồ án I) & 785.14K & 211.80M & 285 & 136 & 10.63 & 6.52$\times$\\
TQNet & 3.39M & 1.22G & 417 & 309 & 25.04 & 15.35$\times$\\
iTransformer & 5.47M & 2.10G & 996 & 747 & 47.30 & 28.99$\times$\\
Informer & 12.45M & 7.91G & 2\,053 & 2\,819 & 176.13 & 107.96$\times$\\
Autoformer & 12.14M & 8.97G & 5\,056 & 5\,579 & 257.63 & 157.91$\times$\\
FEDformer & 18.44M & 8.99G & 3\,295 & 5\,080 & 301.83 & 185.01$\times$\\
TimeMixer & 5.58M & 29.97G & 10\,942 & 25\,557 & 812.01 & 497.73$\times$\\
PatchTST & 8.69M & 16.08G & 10\,949 & 76\,888 & 928.74 & 569.27$\times$\\
TimesNet & 151.32M & 1081.75G & 4\,093 & 147\,660 & 13\,700.21 & 8\,397.57$\times$\\
\bottomrule
\end{tabular}}
\end{table}


Bảng~\ref{tab:efficiency} cho thấy:
\begin{itemize}
  \item DyTimeBase nhanh thứ hai trong 16 mô hình khi suy luận (1.63 ms một mẫu), chỉ sau TimeBase (1.01 ms), và nhanh
        hơn SparseTSF (2.39 ms), FITS, DLinear.
  \item So với các mô hình nhẹ có học chu kỳ, DyTimeBase nhanh hơn CycleNet/Linear 4.1 lần, CycleNet/MLP 4.9 lần và TQNet
        15 lần; TQNet cần 1.22G MACs, gấp khoảng 51 lần DyTimeBase. So với các mô hình Transformer, DyTimeBase nhanh hơn
        iTransformer khoảng 29 lần, TimeMixer khoảng 500 lần và PatchTST khoảng 570 lần.
  \item Chi phí huấn luyện cũng thấp: 40 giây một epoch và 146 MB bộ nhớ, trong khi PatchTST cần gần 11 GB và khoảng
        21 giờ một epoch trên cùng thiết bị.
  \item So với DyTimeNet của Đồ án I ở cùng cấu hình, DyTimeBase có ít hơn khoảng 4.1 lần tham số (190.5K so với 785.1K),
        ít hơn khoảng 8.9 lần MACs (23.8M so với 211.8M) và suy luận nhanh hơn khoảng 6.5 lần.
\end{itemize}
Nguyên nhân là mọi phép biến đổi theo trục thời gian chạy trên các kích thước $N$, $R$, $P$ (vài chục) thay vì $L$, $S$
(hàng trăm), như đã phân tích ở mục~\ref{sec:complexity}.

\subsection{Nghiên cứu loại bỏ (Ablation study)}

Bảng~\ref{tab:ablation} đánh giá đóng góp của từng khối trong nhánh trộn kênh. Các biến thể giữ nguyên mọi siêu tham số
khác, chỉ bỏ đi một khối.

\begin{table}[H]
\centering
\caption{Nghiên cứu loại bỏ (ablation) các khối của nhánh trộn kênh: MSE trung bình trên 4 độ dài dự báo và 3 seed. Dòng cuối giữ nguyên kiến trúc, chỉ tăng tốc độ học của nhánh trộn kênh từ $5\cdot10^{-4}$ lên $2\cdot10^{-3}$. \textbf{Đậm}: tốt nhất mỗi cột.}
\label{tab:ablation}
\small
\resizebox{\textwidth}{!}{\begin{tabular}{@{}lcccccc@{}}
\toprule
Biến thể & ETTh1 & ETTh2 & ETTm1 & ETTm2 & Weather & Trung bình\\
\midrule
DyTimeBase (đầy đủ) & 0.4030 & 0.3580 & \textbf{0.3509} & 0.2634 & 0.2211 & 0.3193\\
\quad bỏ CrossVBlock & 0.4028 & 0.3595 & 0.3554 & 0.2615 & 0.2226 & 0.3204\\
\quad bỏ IntraVBlock & 0.4004 & 0.3668 & 0.3524 & 0.2661 & \textbf{0.2186} & 0.3209\\
\quad bỏ cả nhánh trộn kênh ($=$ TimeBase$^{*}$) & 0.4105 & 0.3627 & 0.3674 & \textbf{0.2539} & 0.2195 & 0.3228\\
\midrule
DyTimeBase, $\alpha_M = 2\cdot10^{-3}$ & \textbf{0.3992} & \textbf{0.3569} & 0.3554 & 0.2627 & 0.2203 & \textbf{0.3189}\\
\bottomrule
\end{tabular}}
\end{table}


\begin{itemize}
  \item Tính trung bình trên 5 bộ dữ liệu, DyTimeBase đầy đủ (0.3193) tốt hơn mọi biến thể đã bỏ khối: bỏ CrossVBlock
        (0.3204), bỏ IntraVBlock (0.3209) và bỏ cả nhánh (0.3228). Cả hai khối đều có đóng góp, và hai khối bổ trợ
        nhau: bỏ một khối thì vẫn giữ được một phần lợi ích.
  \item Trên Weather (21 biến), bỏ CrossVBlock cho kết quả kém nhất (0.2226), còn chỉ giữ CrossVBlock (bỏ IntraVBlock)
        cho kết quả tốt nhất (0.2186), tốt hơn cả TimeBase$^{*}$ (0.2195). Điều này phù hợp với giả thuyết rằng trộn kênh
        có ích hơn khi số biến lớn.
  \item Tăng tốc độ học của nhánh trộn kênh lên $2\cdot10^{-3}$ gần như không thay đổi kết quả (0.3189), nên cấu hình
        $\alpha_M=5\cdot10^{-4}$ của Đồ án I được giữ nguyên.
\end{itemize}

\subsection{Thảo luận và hạn chế}

\begin{itemize}
  \item Các chênh lệch trong nghiên cứu loại bỏ (0.3--1.1\%) cùng cỡ với độ lệch chuẩn giữa các seed, nên chỉ cho thấy xu
        hướng, chưa đủ để khẳng định chắc chắn đóng góp của từng khối.
  \item DyTimeBase chưa vượt được số liệu công bố của TimeBase, và kém các mô hình tuyến tính mạnh trên ETTh2, ETTm2.
  \item Số liệu độ chính xác của các mô hình so sánh (trừ TimeBase$^{*}$, CycleNet và TQNet) được trích dẫn chứ không chạy
        lại. TQNet dùng siêu tham số của script $L=96$ nên có thể chưa tối ưu ở $L=720$. Số tham số của Informer,
        Autoformer, FEDformer ở Bảng~\ref{tab:efficiency} (cấu hình script của Time-Series-Library) khác với số
        trong~\cite{timebase}.
  \item Max Mem đo bằng RAM của CPU; Epoch Time của các mô hình lớn là ước tính từ thời gian một bước huấn luyện. MACs
        không tính phép biến đổi Fourier nhanh, nên FEDformer, FITS và TimesNet bị đếm thiếu một phần nhỏ.
\end{itemize}

\section*{Tóm tắt chương}

DyTimeBase dự báo một mẫu trong 1.63 ms trên CPU, nhanh thứ hai trong 16 mô hình, nhanh hơn CycleNet và TQNet từ 4 đến
15 lần và nhanh hơn các mô hình Transformer từ 29 đến hàng nghìn lần. Về độ
chính xác, DyTimeBase vượt toàn bộ các mô hình Transformer và DLinear, tốt hơn TQNet trên cả 5 bộ dữ liệu, đạt kết quả
tốt nhất trên ETTm1, đứng thứ hai trên ETTh1 và Weather, và ngang ngửa nhóm mô hình tuyến tính mạnh (SparseTSF, FITS,
CycleNet/Linear) cùng PatchTST. So với TimeBase trong cùng quy trình, nhánh trộn kênh cải thiện kết quả trên 3/5 bộ dữ liệu.

%=====================================================================
\begin{thebibliography}{99}
\bibitem{timebase} Q. Huang, Z. Zhou, K. Yang, Z. Yi, X. Wang, Y. Wang, ``TimeBase: The power of minimalism in efficient
  long-term time series forecasting,'' in \emph{Proc. ICML}, PMLR 267, pp. 26227--26246, 2025.
\bibitem{dytimenet} T. Chen, J. Lai, Y. Sun, W. K. Chan, ``DyTimeNet: Dynamic cross-variable dependency network with sparse
  strategy for multivariate time series forecasting,'' \emph{Neurocomputing}, vol. 676, p. 132964, 2026.
\bibitem{tslib} Y. Wang, H. Wu, J. Dong, Y. Liu, M. Long, J. Wang, ``Deep time series models: A comprehensive survey and
  benchmark,'' \emph{arXiv:2407.13278}, 2024.
\bibitem{cyclenet} S. Lin, W. Lin, X. Hu, W. Wu, R. Mo, H. Zhong, ``CycleNet: Enhancing time series forecasting through
  modeling periodic patterns,'' in \emph{Proc. NeurIPS}, 2024.
\bibitem{tqnet} S. Lin, H. Chen, H. Wu, C. Qiu, W. Lin, ``Temporal query network for efficient multivariate time series
  forecasting,'' in \emph{Proc. ICML}, PMLR 267, 2025.
\bibitem{sparsetsf} S. Lin, W. Lin, W. Wu, H. Chen, J. Yang, ``SparseTSF: Modeling long-term time series forecasting with
  1k parameters,'' in \emph{Proc. ICML}, 2024.
\bibitem{timemixer} S. Wang, H. Wu, X. Shi, T. Hu, H. Luo, L. Ma, et al., ``TimeMixer: Decomposable multiscale mixing for
  time series forecasting,'' in \emph{Proc. ICLR}, 2024.
\bibitem{fits} Z. Xu, A. Zeng, Q. Xu, ``FITS: Modeling time series with 10k parameters,'' in \emph{Proc. ICLR}, 2024.
\bibitem{itransformer} Y. Liu, T. Hu, H. Zhang, H. Wu, S. Wang, L. Ma, M. Long, ``iTransformer: Inverted transformers are
  effective for time series forecasting,'' in \emph{Proc. ICLR}, 2024.
\bibitem{dlinear} A. Zeng, M. Chen, L. Zhang, Q. Xu, ``Are transformers effective for time series forecasting?'' in
  \emph{Proc. AAAI}, 2023.
\bibitem{patchtst} Y. Nie, N. H. Nguyen, P. Sinthong, J. Kalagnanam, ``A time series is worth 64 words: Long-term
  forecasting with transformers,'' in \emph{Proc. ICLR}, 2023.
\bibitem{timesnet} H. Wu, T. Hu, Y. Liu, H. Zhou, J. Wang, M. Long, ``TimesNet: Temporal 2D-variation modeling for general
  time series analysis,'' in \emph{Proc. ICLR}, 2023.
\bibitem{fedformer} T. Zhou, Z. Ma, Q. Wen, X. Wang, L. Sun, R. Jin, ``FEDformer: Frequency enhanced decomposed transformer
  for long-term series forecasting,'' in \emph{Proc. ICML}, 2022.
\bibitem{autoformer} H. Wu, J. Xu, J. Wang, M. Long, ``Autoformer: Decomposition transformers with auto-correlation for
  long-term series forecasting,'' in \emph{Proc. NeurIPS}, 2021.
\bibitem{informer} H. Zhou, S. Zhang, J. Peng, S. Zhang, J. Li, H. Xiong, W. Zhang, ``Informer: Beyond efficient
  transformer for long sequence time-series forecasting,'' in \emph{Proc. AAAI}, 2021.
\end{thebibliography}

\end{document}
